Legacy Java systems pose a long-term challenge for developers: they typically lack documentation and automated tests, which makes every change risky.
 Introducing unit tests would reduce this risk, but writing them retroactively is time-consuming and often hindered by poor testability, especially tight dependencies between components. 
 The aim of this thesis is to explore the use of large language models (LLMs) for automated unit test generation for such systems.
  We propose a methodology for generating and evaluating tests and practically assess the capabilities of several LLMs. 
  The models are compared empirically in terms of generation success rate, code coverage and test reliability.
   The thesis also analyses how structural properties of source code, such as cyclomatic and cognitive complexity and the degree of dependencies, affect testability and the success of test generation. 
   In addition, we examine whether AI can refactor code to make it more testable.
%  The results show that [fill in main findings]. The thesis contributes [fill in contribution, e.g. practical recommendations for applying LLMs to testing legacy code].